\documentclass[11pt]{article}
\usepackage[T1]{fontenc}
\usepackage{amsmath,amssymb,amsthm}
\usepackage[margin=1in]{geometry}
\usepackage[hidelinks]{hyperref}
\newtheorem{theorem}{Theorem}
\newtheorem{lemma}[theorem]{Lemma}
\newcommand{\DD}{d'd''}
\newcommand{\MA}{\operatorname{MA}}
\newcommand{\C}{\mathcal C}
\newcommand{\G}{\mathcal G}
\title{Two distinct residue wells carrying the entire raw mass}
\author{Local unverified proof draft}
\date{9 October 2026}
\begin{document}
\maketitle
\begin{abstract}
On the supplied analytic quotient, two centres with distinct residues in
both coordinates admit a genuine continuous CLD weight whose actual
Monge--Ampere measure is one atom in each residue bidisc and no Gauss atom.
An exact cross-product calculation extends this to two independent Radon
depth distributions, each of mass at most one. Their full dual is finite
exactly when the sum of their tail-to-the-power-$3/2$ integrals is finite.
Actual continuous decreasing presentations, including unbounded examples
with the entire mass off the canonical skeleton, are constructed.
These are special prescribed-support equations. General recognition,
the single whole-mass Dirac equation, genuine nonalgebraization/descent,
and the complete independent incoming-proof audit remain unresolved.
\end{abstract}

\section*{Domain, definitions, and external boundary}
Retain precisely the supplied compact, good, strict, smooth, separated,
boundaryless pure surface
\[
 K=\mathbb C((s)),\quad |s|=e^{-1/8},\quad t=s^8,
 \qquad X=(\mathbb G_m^2)^{\rm an}/\langle(t,2),(3,t)\rangle.
\]
Write $r_i=-\log|z_i|$, $\rho=(r_1^2+r_2^2)/2$, and $\theta=\DD\rho$.
Deck transformations translate $r$ by integers, so reference changes are
smooth harmonic affine functions. The class $\C$ is the ENTIRE original
minimal continuous CLD sheaf of weights, not a finite-log-max subclass:
$\rho+a$ locally belongs to the smallest continuous sheaf containing
smooth positive-Hessian functions and closed under local uniform limits.
Set $\theta_a=\theta+\DD a$, $B(a,b)=\theta_a\wedge\theta_b$,
\[
 E(a)=\tfrac13\int_Xa(\theta^2+B(a,0)+\MA(a)),\qquad
 D(\mu)=\sup_{a\in\C}\bigl(E(a)-\int_Xa\,d\mu\bigr).
\]
The reference raw mass is two. The domain $\G$ requires an ACTUAL
continuous decreasing $\C$ presentation with finite lower energy.
The quotient's geometry is supplied; its genuine nonalgebraization and
coherent descent are not independently proved here.

The external results used below are from Chambert--Loir--Ducros,
\emph{Formes differentielles reelles et courants sur les espaces de
Berkovich}, version 2, 25 July 2025 (CLD v2). Their immediate supplied
statements and proofs were read, with ranges in the separate ledger.
Theorem 15.7.1: a topologically separated pure-dimensional analytic space,
with its actual analytic boundary, has the compatible calibrated tropical
structure. We apply it to $X$ and ordinary open subsets, all boundaryless.
Definition 8.1.7/Lemma 8.1.8 give exactly the minimal continuous sheaf
above, including every ordinal stage. Corollary 8.2.4 permits convex
coordinatewise increasing continuous composition.
Lemmas 8.5.1 and 8.5.5 and Propositions 8.7.2, 8.7.4 give positive
recursive products, local uniform continuity and symmetry for continuous
psh potentials and positive closed currents on boundaryless tropical spaces.
Section 8.8.5 gives their multilinear dpsh extension and invariance under
smooth $\DD$-null changes. Proposition 8.5.6 gives
\[
 \DD u\wedge j_*S=j_*(\DD(j^*u)\wedge S)
\]
for a topologically proper tropical morphism $j$ and a positive closed
current $S$, used only on fixed closed analytic fibres. Propositions
7.2.8--13 identify positive top currents with Radon measures and upgrade
current convergence to continuous tests. Proposition 5.11.1 is smooth
Green/Stokes with the actual boundary term, zero on our compact $X$.
Sections 14.5.2--3 calibrate an identity torus moment by its analytic
degree and saturated lattice; here both are one. Proposition 5.10.1(c),(d)
and 5.10.2 give calibrated integration and Radon reconstruction.
Corollaries 6.2.4--6 give compactly supported smooth cutoffs and uniform
smooth approximation on locally differentially separated affine tropical
spaces. Proposition 6.2.7 gives smooth partitions for a paracompact such
space; Corollary 6.2.8 gives nonnegative, compactly supported, locally
finite cutoffs summing to one near a compact set. Our compact Hausdorff
$X$ is paracompact and its ordinary interior is locally differentially
separated by Theorem 15.7.1. These statements and their complete immediate
proofs were additionally read in the full CLD carrier; their Appendix A
proofs remain the disclosed conventional topological boundary.

The decisive analytic formula is Corollary 18.6.2: for analytic
$f_1,\ldots,f_d$ on a boundaryless, topologically separated,
equidimensional analytic space and positive real $t,R_i$,
\begin{equation}\label{fibre}
 (\DD\log\max(t,|f_1|/R_1,\ldots,|f_d|/R_d))^d
 =\delta_{f^{-1}(\eta_{(tR_1,\ldots,tR_d)})}.
\end{equation}
The right side is the actual completed-residue-field fibre integration
current. Proposition 18.6.1 proves the generically-flat relative formula
used in its proof; the residual projective intersection coefficient there
is one. Literal prose slips in 18.6.2 are not used to assert vanishing of
the Hessian: the dimension argument concerns its TOP POWER. Our
translations, inversions and primitive monomial changes are explicit
isomorphisms. No field-extension multiplicity or value-group restriction
is added. These source applications keep the conventional smooth
partitions/density, Berkovich boundary, and relative-intersection inputs
of CLD. This is not an audit of their entire recursive foundations.
No previous note or verdict is accepted as a correctness grant.

\section*{Whole-class bounded calculus needed for the full dual}
For clarity the full-$\C$ tangent input is reconstructed, rather than
assumed from a previous note. For each fixed positive closed $(1,1)$
current $T$, complete compact smooth gradients in the positive pairing
$\int d'f\wedge d''g\wedge T$, with nonnegative local cutoffs.
This gradient is closable under uniform convergence: if $f_j\to0$
uniformly and their gradients converge strongly, pairing against a compact
smooth one-form $\beta$ gives
$\int d'f_j\wedge J\beta\wedge T=-\int f_jd'(J\beta)\wedge T\to0$.
The coefficient measures are finite on this support and smooth tests are
dense in the completed space. Smooth cutoffs and Leibniz patch the domains.
If $L_Tf=\DD(fT)$ is signed Radon, integration against smooth tests,
then uniform and strong approximation of the other factor, gives Green
\[
 Q_T(f,g)=-\int gL_Tf=-\int fL_Tg,\qquad Q_T(f)\ge0.
\]
The smooth chain rule passes strongly for functions with bounded continuous
derivative. With nested smooth cutoffs $\chi,\eta$, $\eta=1$ near
$\operatorname{supp}\chi$, the same identity gives
\[
 Q_{T,\chi}(f)=\tfrac12\int f^2\DD\chi\wedge T-\int\chi fL_Tf.
\]
For uniformly convergent psh-domain functions $f_j$, bounded by $A$ on
$\operatorname{supp}\eta$, put $\epsilon=\|f_j-f_k\|$ there.
Positivity of $L_Tf_j$ and
$\int\eta L_Tf_j=\int f_j\DD\eta\wedge T$ imply
\[
 Q_{T,\chi}(f_j-f_k)\le
 \tfrac12\epsilon^2\|\DD\chi\wedge T\|_{\rm TV}
 +2A\epsilon\|\chi\|_\infty\|\DD\eta\wedge T\|_{\rm TV}.
\]
Thus their gradients are Cauchy, and choosing smooth approximants with
both errors tending to zero proves membership of the limit. Distributional
positivity persists. This uniformly closed sheaf contains all smooth
positive-Hessian generators, hence the ENTIRE minimal sheaf. Differences
and a finite partition prove global Green for all bounded weights.
This argument uses no arbitrary Borel multiplier or CLD 8.6.7.

Symmetric product expansion and Green now give, for
$a_h=(1-h)a+hb$,
\begin{equation}\label{calculus}
 \frac{d}{dh}E(a_h)=\int(b-a)\MA(a_h),\qquad
 \frac{d^2}{dh^2}E(a_h)=-2Q_{\theta_{a_h}}(b-a)\le0.
\end{equation}
Indeed differentiating the cubic definition and moving a $\DD$ factor
by Green produces three equal first-variation terms; differentiating once
more gives $2\int(b-a)\DD(b-a)\wedge\theta_{a_h}=-2Q$.
All bounded mixed masses are two: their difference from $\theta^2$ is
a sum of globally $\DD$-exact currents, whose value on the global test
$1$ is zero. We also obtain monotonicity, $E(a+k)=E(a)+2k$, and the
cubic cocycle
\[
 E(b)-E(a)=\tfrac13\int(b-a)(\MA(a)+B(a,b)+\MA(b)).
\]
In particular, for every actual bounded $q\in\C$,
\begin{equation}\label{dual}
 E(a)-\int a\MA(q)\le E(q)-\int q\MA(q)\quad(a\in\C).
\end{equation}
This is an inequality on the entire original class; equality at $a=q$
computes its full dual.

\section*{Baseline and actual two-residue log maximum}
Put $\xi_i=r_i-\lfloor r_i\rfloor$ and
\[
 q_0=\sum_{i=1}^2\xi_i(1-\xi_i)/2-1/4,
 \qquad g=\sigma(0).
\]
On the injective ordinary chart $U=\{|r_i|<1/4\}$, with
$A_i=\max(1,|z_i|)$, its local potential is
\[
 \rho+q_0=\log(A_1A_2)-\tfrac12\log|z_1z_2|-1/4.
\]
At integer walls the one-coordinate local potential has a positive
slope jump; elsewhere it is affine after adding $\rho$. Smooth convex
maxima establish $q_0\in\C$ on all quotient charts. Its square is
$2\delta_g$: each coordinate Hessian is the fixed unit Gauss curve by
\eqref{fibre}, its own square vanishes by one-coordinate smooth
approximation, and their mixed product restricts to the other coordinate
Gauss point, coefficient one. Outside the simultaneous integer walls the
square is zero. Identity calibration also gives
$\theta^2=2\sigma_*dr_1dr_2$ and
$B(q_0,0)=\sigma_*(\delta_{r_1=0}dr_2+dr_1\delta_{r_2=0})$.
For the first assertion, off the joint Gauss section a Laurent polynomial
has norm strictly below its monomial maximum. This persists ordinarily
and forces coordinate tropicalization into its finite hypersurface of
dimension at most one, where the coordinate top form vanishes.
On the section the identity degree-one calibration gives the raw factor
$2!$. The curve formula similarly gives the second assertion. Therefore
\[
 \int q_0\theta^2=-1/6,\quad
 \int q_0B(q_0,0)=-1/3,\quad
 \int q_0\MA(q_0)=-1/2,\quad E(q_0)=-1/3.
\]
These integrals follow from $\int_0^1x(1-x)/2\,dx=1/12$.

Choose units $a_1,a_2,b_1,b_2\in K$ with
$|a_1-a_2|=|b_1-b_2|=1$. They may concretely be $a_1=b_1=1$,
$a_2=b_2=2$. Let
\[
 D_i=\{|z_1-a_i|<1,\ |z_2-b_i|<1\}\subset U,
 \qquad P_i=(a_i,b_i).
\]
These are disjoint ordinary open bidiscs. Let $x_{i,c}$ be their joint
Gauss point of radius $e^{-c}$, $c>0$, and set $x_{i,0}=g$.
Define $\gamma_i=\max(0,-\log\max(|z_1-a_i|,|z_2-b_i|))$ on $U$,
zero elsewhere, and $\gamma_i(P_i)=+\infty$.
It is extended continuous, supported on $D_i$, and
$\gamma_i(x_{j,c})=c$ if $i=j$ and zero otherwise.
Put $v_{i,c}=-\min(c,\gamma_i)$ and $q_{i,c}=q_0+v_{i,c}$.

\begin{theorem}\label{two}
For every real $c_1,c_2>0$, the weight
$q_{12}=q_0+v_{1,c_1}+v_{2,c_2}$ belongs to the original $\C$, and
\[
 \MA(q_{12})=\delta_{x_{1,c_1}}+\delta_{x_{2,c_2}}.
\]
Thus its entire raw mass two is nonskeletal. It has a bounded actual
decreasing presentation $q_{12}+2^{-n}$ in $\G$.
\end{theorem}
\begin{proof}
Set $u=z_1-a_1$, $v=z_1-a_2$, $p=z_2-b_1$, $q=z_2-b_2$ and
$R_i=e^{-c_i}<1$. At EVERY analytic point, the strong triangle inequality
and the unit differences imply
$\max(|u|,|v|)=A_1$ and $\max(|p|,|q|)=A_2$.
If $A_1>1$, both $|u|,|v|$ equal $A_1$; likewise for $A_2>1$.
Consider the actual polynomial log maximum
\[
 H=\log\max(|uq|,|vp|,R_1|v|,R_1|q|,R_2|u|,R_2|p|).
\]
On $D_1$, $|v|=|q|=1$, so $H=\log\max(R_1,|u|,|p|)$.
On $D_2$, $|u|=|p|=1$, so $H=\log\max(R_2,|v|,|q|)$.
Outside both discs, $\max(|uq|,|vp|)=A_1A_2$.
To verify the last assertion, if either $A_i>1$, the corresponding two
norms are equal and the other pair has maximum $A_j$. If both $A_i=1$
and $|u|<1$, then $|v|=1$ and being outside $D_1$ forces $|p|=1$;
if $|v|<1$ use $|u|=|q|=1$. If $|u|=|v|=1$, one of $|p|,|q|$ is one.
Each cap is at most $A_1A_2$ since $A_i\ge1$ and $R_i<1$.
This proves the identity, including zeros and nonclassical points,
\[
 H=\log(A_1A_2)+v_{1,c_1}+v_{2,c_2}.
\]
The maximum has no common zero. Indeed the common zeros of $uq,vp$
are precisely $P_1,P_2$, and at these the corresponding opposite-coordinate
caps have norm $R_i>0$. Near any point, slots vanishing there can be
omitted since a nonzero slot has a positive local lower bound. The remaining
slots are unit logarithms. Their smooth convex maxima have positive
Hessian and converge locally uniformly. Subtracting
$\tfrac12\log|z_1z_2|+1/4$ is smooth harmonic. Thus $\rho+q_{12}$
is genuinely continuous psh on $U$.
If a $|z_j|$ differs from one, then for each centre
$|z_j-a_i|\ge1$ (or the corresponding $b_i$), so both corrections vanish.
In particular they vanish on the ordinary collar
$1/8<\max|r_i|<1/4$. This glues to $q_0$ on the quotient. The translated
coordinates were used only on the injective chart; deck units $2,3$ were
not assumed to fix them.

On $D_i$, the local potential differs smoothly harmonically from
$\log\max(R_i,|z_1-a_i|,|z_2-b_i|)$. Apply \eqref{fibre} with
$d=2,t=R_i,R_1=R_2=1$ to the translation isomorphism of this ordinary
open bidisc. Its fibre over the joint radius-$R_i$ Gauss point is the
single point $x_{i,c_i}$, degree one. All source hypotheses hold on
$D_i$, and $R_i$ may be outside $|K^\times|$.
Hence each disc carries exactly one unit of the actual positive Radon
measure. The total mass is two by the global exact-current argument.
Positivity leaves no mass on the complement: the difference from these
two atoms is positive with total mass zero. This proves equality on all
analytic continuous tests, with no inverse projected-measure convention.
Additionally, near $g$ the caps are strictly inactive, $uq,vp$ are units,
and $H=\log|vp|+\max(0,\log|uq/vp|)$ has rank-one smooth approximations;
its actual square is zero there. Boundedness and the stated decreasing
presentation are immediate.
\end{proof}

\section*{Cross-products and independent masses}
The one-well weights also belong to $\C$: on $U$ their potential is
\[
 \log\max(e^{-c},|z_1-a_i|,|z_2-b_i|,
                 |(z_1-a_i)(z_2-b_i)|)
 -\tfrac12\log|z_1z_2|-1/4.
\]
The positive constant and the same ordinary collar prove this at zeros
as well. Formula \eqref{fibre} gives coefficient one at $x_{i,c}$.
Near $g$, subtract the unit log of the product and use inversion of
the two translated coordinates; its unit joint Gauss fibre is $g$,
also coefficient one. Positivity and mass two yield
\[
 \MA(q_{i,c})=\delta_g+\delta_{x_{i,c}}.
\]
Moreover $B(q_0,q_{i,c})=2\delta_g$. To see this on actual tests, restrict
to either fixed original-coordinate unit Gauss fibre. Its fixed coordinate
has $|z_j-a_i|=1$ (respectively $|z_j-b_i|=1$), and the above logarithmic
maximum restricts to $\log\max(1,|z_k|)$. The proper fixed-fibre
restriction formula and the curve version of \eqref{fibre} give $g$,
coefficient one, for each of the two fibres.

Writing $\Delta_{i,c}=\DD v_{i,c}$, we consequently have
$\theta_{q_0}\Delta_{i,c}=0$ and
$\Delta_{i,c}^2=\delta_{x_{i,c}}-\delta_g$. Theorem \ref{two},
for every pair of depths, and multilinearity then imply the crucial
ACTUAL cross-product identity
\begin{equation}\label{cross}
 \Delta_{1,c}\wedge\Delta_{2,d}=0\qquad(c,d>0).
\end{equation}
This was derived from positive global weights and fully identified actual
measures; it was not inferred merely from disjoint negative wells.

For completeness the same-branch product is
\begin{equation}\label{same}
 B(q_{i,c},q_{i,d})=\delta_g+\delta_{x_{i,\min(c,d)}}.
\end{equation}
Here is its independent proof. Take $c<d$, omit polynomial zero slots,
and put $b_j=\log|z_j-\text{centre}_j|$, $m=\max(b_1,b_2)$.
Where some $b_j\ge0$ the two maxima locally agree, and their square is
the already identified Gauss atom. In the negative quadrant, if $m<-c$
the $c$-maximum is constant; if $m>-c$ both equal $m$, whose square
is rank one. At $m=-c,b_1\ne b_2$ the $d$-maximum is one unit log.
At $b_1=b_2=-c$, its Hessian is
$\DD\log\max(1,|w_1/w_2|)$, the fixed ratio Gauss curve by
\eqref{fibre} with $d=1$. On this curve over $\mathscr H(\eta_1)$,
the other maximum restricts to $\log\max(e^{-c},|w_2|)$.
Proper fixed-fibre restriction and the curve formula give one point.
It is $x_{i,c}$: the primitive exponent bijection
$(j,k)\mapsto(j,j+k)$ identifies the iterated Gauss norm in
$(w_1/w_2,w_2)$ with the joint norm in $(w_1,w_2)$, with determinant one.
At polynomial zeros both potentials locally involve only the remaining
coordinate, so the mixed square vanishes. The equal-depth case was proved
above. This covers all analytic points and proves \eqref{same}.

\begin{theorem}\label{independent}
For arbitrary $0\le\lambda_i\le1$ and depths $c_i>0$, put
$q=q_0+\lambda_1v_{1,c_1}+\lambda_2v_{2,c_2}$. Then $q\in\C$ and
\[
 \MA(q)=(2-\lambda_1^2-\lambda_2^2)\delta_g
        +\lambda_1^2\delta_{x_{1,c_1}}+\lambda_2^2\delta_{x_{2,c_2}},
\]
\[
 E(q)=-\tfrac13-\tfrac13(\lambda_1^3c_1+\lambda_2^3c_2),\qquad
 D(\MA(q))=\tfrac16+\tfrac23(\lambda_1^3c_1+\lambda_2^3c_2).
\]
Each off-Gauss mass may independently reach one.
\end{theorem}
\begin{proof}
The weight is the convex combination of the four actual corners
$q_0,q_{1,c_1},q_{2,c_2},q_{12}$ with respective coefficients
$(1-\lambda_1)(1-\lambda_2),\lambda_1(1-\lambda_2),
(1-\lambda_1)\lambda_2,\lambda_1\lambda_2$. Thus positivity has
already been established, even when $\lambda_1+\lambda_2>1$.
Equations \eqref{cross} and the single-well products give the measure.
Along $q_0+h(q-q_0)$ the off-Gauss coefficients are $h^2\lambda_i^2$,
and $q-q_0$ has value $-\lambda_ic_i$ at $x_{i,c_i}$ and zero at $g$.
Equation \eqref{calculus} gives energy derivative
$-h^2\sum\lambda_i^3c_i$, hence the energy formula.
The target moment is $-1/2-\sum\lambda_i^3c_i$ since $q_0=-1/4$
at all target points. Equation \eqref{dual} computes the entire dual.
The normalized weight is nonpositive and has supremum zero, achieved at
a maximum of $q_0$ outside $U$.
\end{proof}

\section*{Two arbitrary Radon depth distributions, with full nonskeletal mass}
Let $\zeta_i$ be any finite positive Radon measure on $(0,\infty)$,
$M_i=\zeta_i((0,\infty))\le1$. Define
\[
 T_i(t)=\zeta_i((t,\infty)),\quad
 d_i(c)=\int_0^c\sqrt{T_i(t)}dt,\quad
 S_i=\int_0^\infty T_i(t)^{3/2}dt,
\]
and the actual target
\[
 \mu=(2-M_1-M_2)\delta_g+\sum_{i=1}^2(c\mapsto x_{i,c})_*\zeta_i.
\]
The branch maps are continuous, including at depth zero: polynomial Gauss
norms are finite maxima of continuous monomial norms, and Laurent norms
follow since the coordinates are units. Compact depth restrictions give
Radon pushforwards, and total-variation approximation by those restrictions
gives the full finite Radon measure. Its mass is two and it has no mass
at either infinite-depth rigid centre.

\begin{theorem}\label{radon}
For precisely these targets,
\[
 D(\mu)<\infty\quad\Longleftrightarrow\quad S_1+S_2<\infty.
\]
If this holds, $q=q_0-d_1(\gamma_1)-d_2(\gamma_2)$ has an actual
continuous decreasing original-$\C$ presentation of finite lower energy,
and its canonical presentation-generated operator satisfies
\[
 \MA(q)=\mu,\quad E(q)=-\tfrac13-\tfrac{S_1+S_2}3,
 \quad\int q\,d\mu=-\tfrac12-S_1-S_2,
 \quad D(\mu)=\tfrac16+\tfrac23(S_1+S_2).
\]
It is bounded continuous iff both $\int_0^\infty\sqrt{T_i(t)}dt$ are
finite. Otherwise it is extended continuous, finite away from the centres
with divergent integrals, where its value is $-\infty$.
When $M_1=M_2=1$, the ENTIRE measure is nonskeletal, including in the
unbounded finite-energy cases.
\end{theorem}
\begin{proof}
First take finite atomic depths $c_{i,1}<\cdots<c_{i,N_i}$ with masses
$m_{i,k}$. Set
$\alpha_{i,k}=\sqrt{\sum_{j\ge k}m_{i,j}}$,
$\beta_{i,k}=\alpha_{i,k}-\alpha_{i,k+1}$, and $\alpha_{i,N_i+1}=0$.
Use weights $q_0+\sum_{i,k}\beta_{i,k}v_{i,c_{i,k}}$.
For each branch, its $\beta$ coefficients and $1-\sqrt{M_i}$ form a
probability on its depth slots and the zero correction. Tensoring the
two probabilities expresses the weight as a convex combination of the
actual two-well corners in Theorem \ref{two}. Thus it lies in $\C$.
Equations \eqref{same} and \eqref{cross} give at each depth the coefficient
$\beta_{i,k}^2+2\beta_{i,k}\sum_{j>k}\beta_{i,j}
=\alpha_{i,k}^2-\alpha_{i,k+1}^2=m_{i,k}$.
The remaining coefficient is $2-M_1-M_2$ at $g$.
Its branch correction equals $-d_i(\gamma_i)$ by differentiation between
the finitely many depths and value zero at depth zero. Tonelli gives
$\int d_i(c)\zeta_i(dc)=\int T_i(t)^{3/2}dt=S_i$;
atoms at threshold endpoints change no Lebesgue integral. The same actual
first-variation calculation as above gives $E=-1/3-\sum S_i/3$.

For finite $C>0$, push depths to $\min(c,C)$, keeping the same masses,
and write $d_{i,C}(c)=\int_0^{\min(c,C)}\sqrt{T_i(t)}dt$ and
$S_{i,C}=\int_0^C T_i(t)^{3/2}dt$.
Approximate each depth measure by partition masses placed at positive
right endpoints. Tails converge at every non-atomic threshold, hence
almost everywhere; their square roots are bounded by one. Dominated
convergence gives $L^1(0,C)$ convergence of the square-root tails.
The corresponding weights converge uniformly, with error at most the
sum of these $L^1$ errors, to
$q_C=q_0-d_{1,C}(\gamma_1)-d_{2,C}(\gamma_2)$.
The ENTIRE minimal sheaf and actual recursive-product continuity establish
$q_C\in\C$ and its actual measure
\[
 \MA(q_C)=(2-M_1-M_2)\delta_g
       +\sum_i(c\mapsto x_{i,\min(c,C)})_*\zeta_i.
\]
Uniform energy continuity or first variation gives
$E(q_C)=-1/3-\sum_iS_{i,C}/3$. Tonelli gives both its own target moment
and its moment against the untruncated $\mu$ as $-1/2-\sum_iS_{i,C}$.

The actual $q_n$ are continuous, normalized, nonpositive and decreasing
on EVERY point of $X$. If $\sum S_i<\infty$, their energies have a finite
lower bound, so they are a required $\G$ presentation. For any continuous
analytic test $h$, the difference of its integrals against $\MA(q_n)$
and $\mu$ is at most $2\|h\|_\infty\sum_i\zeta_i((n,\infty))$.
Thus the actual bounded operators converge weakly to $\mu$.
The canonical generated-operator lemma proved next identifies this as
$\MA(q)$; no Borel limit is recognized by declaration.
Tonelli proves the limiting target moment. Applying \eqref{dual} at
$q_n$ to each fixed $a\in\C$ and passing weakly gives the asserted upper
bound for FULL $D(\mu)$. The actual competitors $q_n$ give the reverse
bound, since their objectives are exactly $1/6+2\sum S_{i,n}/3$.
If $\sum S_i=\infty$, these same bounded competitors give $D=\infty$;
no operator is assigned to that ineligible limit. Boundedness and extended
continuity follow from the actual extended depth functions and their
continuous increasing integrals. The two ordinary wells are disjoint,
so the other correction vanishes in a neighbourhood of either centre.
\end{proof}

\section*{Canonicity only for actual generated presentations}
Here is the promised proof of the operator fact; it makes no assertion
about an unrepresented upper relaxation. The real span $V$ of $\C$ is
dense in $C(X)$. To verify this without a skeletal replacement, $V$ is a
vector lattice: write its elements as $A(p-q)$ with $p,q\in\C$, filling
convex averages with zero to obtain a common $A$; then
\[
 \max(A(p-q),A(r-s))=2A[\max((p+s)/2,(r+q)/2)-(q+s)/2].
\]
It contains constants and separates actual points. Different tropical
values are separated by small ordinary tropical-ball quadratic caps.
At equal tropical values, choose lifts in one injective chart and a
Laurent polynomial with different seminorms; some
$\psi=\log\max(1,|P|/R)$ has different values. It is bounded continuous
psh on a compact surrounding polyannulus, including at zeros by a smooth
convex ramp constant sufficiently far to the left. Centre the reference
quadratic $\rho_0$ at the common tropical value. Choose $\delta>0$ and
$\epsilon>0$ small so
$w=\max(0,\delta+\epsilon\psi-\rho_0)$ is zero on an outer ordinary
collar, and positive at the two central points. Then $\rho_0+w$ is a
maximum of two actual psh slots, glues globally, and separates the points.
For a continuous $h$ and tolerance $\eta$, separation and constants give
an element of $V$ matching $h$ at any pair $x,y$. For fixed $x$, finitely
many neighbourhoods of $y$ give a finite maximum above $h-\eta$
everywhere and equal to $h$ at $x$. Finitely many neighbourhoods of $x$
then give a finite minimum uniformly within $\eta$ of $h$. This proves
density on compact Hausdorff $X$ without countability assumptions.

Let $b_n\downarrow v$ be any ACTUAL presentation with finite lower energy.
Monotonicity makes $e=\lim E(b_n)$ finite. Given another presentation
$c_m$ and fixed $n,\epsilon>0$, the open sets
$\{c_m<b_n+\epsilon\}$ increase and cover compact $X$; a finite cover
gives a single $m$ with the inequality everywhere. Monotonicity and
$E(b_n+\epsilon)=E(b_n)+2\epsilon$ prove independence of $e$ by comparison
in both directions. The same holds for finite convex combinations,
whose lower energy bounds follow from concavity.
For fixed bounded $w\in\C$, the cubics
$P_n(t)=E((1-t)b_n+tw)$ and $A_n(t)=E(tb_n)$ converge at every $t\in[0,1]$.
Four distinct interpolation parameters imply coefficient convergence,
and their limits are independent of the actual presentation.
First variation gives
$\int w\MA(b_n)=P'_n(0)+A'_n(1)$, hence a canonical limit.
The dense span $V$, positivity and common mass two extend these limits to
all continuous tests, with error at most $2\|h-w\|_\infty$ for a single
measure. The limiting positive norm-two functional has, by the
compact-Hausdorff Riesz theorem, a unique positive Radon measure of mass
two. This is the canonical presentation-generated operator, obtained
from genuine bounded products. It applies in Theorem \ref{radon} because
its displayed $q_n$ meet the actual antecedent.

\section*{Explicit full-mass unbounded example and remaining gates}
Take on EACH branch $\zeta_i=c^{-2}dc$ on $[1,\infty)$, of mass one.
Its tail is one below depth one and $t^{-1}$ above it. Thus $S_i=3$,
while $\int\sqrt{T_i}=\infty$. The target
\[
 \mu=\int_1^\infty(\delta_{x_{1,c}}+\delta_{x_{2,c}})c^{-2}dc
\]
has no canonical-skeleton mass at all, and the actual generated solution
is $q_0-d(\gamma_1)-d(\gamma_2)$, where $d(c)=c$ up to one and
$d(c)=2\sqrt c-1$ thereafter. Its only poles are $P_1,P_2$,
and $E=-7/3$, $D=25/6$, $\int q\mu=-13/2$.
For $n\ge1$, $E(q_n)=-7/3+4/(3\sqrt n)$.
For the bounded depths $c_1=c_2=1/8$ and masses one each,
$E(q_{12})=-5/12$ and $D=1/3$.

\begin{lemma}\label{no-coalescence}
Allow the two unit centres and all their eligible depth distributions to
vary within the distinct-residues-in-both-coordinates family just proved.
Every weak Radon limit of their measures assigns mass at most one to
any fixed ordinary residue bidisc $D=\{|z_1-a|<1,|z_2-b|<1\}\subset U$.
Consequently this family and its weak limits cannot give
$2\delta_{x_{a,b,c}}$ for $c>0$.
\end{lemma}
\begin{proof}
A positive-depth Gauss point belongs to $D$ exactly when its centre has
the same residue as $(a,b)$ in both coordinates. At most one of the two
centres can have those residues. Its branch mass is at most one, and
$g\notin D$. Thus every measure in the family has mass at most one on
$D$. For any continuous $0\le h\le1$ compactly supported in $D$,
weak convergence gives $\int h\,d\lambda\le1$. Regularity and the
compact-Hausdorff cutoff characterization of open-set mass give
$\lambda(D)\le1$. But $x_{a,b,c}\in D$, and the whole Dirac has
mass two there. This is solely a restriction on the newly constructed
measure family, not on the entire original $\C$ or $\G$.
\end{proof}

This independently constructs a larger, full nonskeletal mass equation
sector, rather than repeating the earlier one-branch mass budget.
The potentials use polynomial residue data; they are not single-centre
radial fans with the earlier outside-quadrant boundary values.
It does not produce $2\delta_{x_{1,1/8}}$: distinct residues in both
coordinates are essential in the log-max identity, and merging the
centres invalidates that identity. No general dominated analytic
recognition, arbitrary-support existence, or new full-$\G$ uniqueness
audit is asserted. No compactness/envelope replacement for every eligible
measure has been constructed. The entire incoming standalone was retained
but not entirely physically read or independently audited in this run;
the same is true of its historical companion. The exact original residual
and task are retained byte unchanged. P16.7 remains open; candidate is
false. Compilation and independent correctness remain unverified.
\end{document}
